\documentclass[a4paper,12pt,ja=standard,lualatex]{bxjsarticle}
\usepackage{luatexja} 
\usepackage{graphicx}
\usepackage{url}
\usepackage{listings}
\usepackage{makecell}
\usepackage{amsmath, amssymb}
\usepackage{tikz}
\usetikzlibrary{shapes.geometric, arrows.meta, positioning}

\begin{document}
\pagestyle{empty}

\begin{center}
\vspace*{1cm}
\large
{\LARGE 卒業研究報告書}\\
\vspace*{0.8cm}
題目\\
\vspace*{1cm}
{\Huge{音声認識を用いた作業ログ}}\\
{\Huge{自動記録システムの構築と}}\\
{\Huge{業務効率化に関する研究}}\\
\vspace{3mm}

\vspace*{3cm}
指導教員\\
\vspace*{0.3cm}
{\LARGE 森山 真光 准教授}\\ 
\vspace*{3cm}
報告者\\
\vspace*{0.3cm}

{23-1-211-0159}\\
\vspace*{0.3cm}
{\Huge 内山 敦也}\\
\vspace*{0.5cm}
近畿大学情報学部情報学科\\
\vspace*{2cm}
令和8年9月30日提出\\
\end{center}

\newpage
\normalsize

\clearpage
\begin{center}
{\bf \Large 概要}
\end{center} 
本研究では、音声認識技術および近距離無線通信を活用して現場作業内容および位置・作業状況をリアルタイムに記録し、業務効率化を図るハードウェア・ソフトウェア一体型システムの構築を目的とする。
従来の作業記録は手入力や紙媒体に依存しており、作業中断に伴う手間や記録漏れ、作業位置・状況の把握漏れといった課題が存在する。
こうした背景のもと、本研究では「リソースが制限された小型エッジマイコン環境において、認識対象を『加工開始』『加工完了』『段取開始』『段取終了』の4種類の特定作業フレーズに限定し、物理空間情報（BLE電波）と言語解析（Whisper/GPT）を組み合わせることで、現場騒音下での誤認識・誤検知をどこまで抑制できるか」という学術的な問いを設定する。
本研究では、M5CoreS3マイコンデバイスを用いた音声収集・BLEスキャン処理と、サーバ側での高精度な音声認識（Whisper）および目的フレーズへの意図判別・構造化（GPT）を組み合わせたエッジ・クラウド協調型システムを提案する。
発話内容からの正確な作業ログ分類および自動属性付与を実現し、製造現場や点検業務等における入力コストの削減と作業プロセスの透明化・効率化への貢献を目指す。

\newpage

\tableofcontents
\newpage
\setcounter{page}{1}
\pagestyle{plain}

\section{序論}\label{sec1}
本研究では、音声認識技術およびIoTデバイスを用いて作業内容や現場情報をリアルタイムに記録し、業務効率化を図るシステムの構築を目的とする。

\subsection{本研究の背景}
製造業や設備点検をはじめとする多くの作業現場において、作業ログの記録は日報の手入力やチェックシートへの記入などアナログな手段に依存している。
こうした入力作業は本来の業務を停止させるため、生産性の低下を招くだけでなく、手入力に伴う入力ミスや記録漏れ、実作業時間の不正確さといった課題を引き起こす。
特に工具や機械を扱う作業においては、作業の途中で記録端末を操作することが物理的に困難な場合が多く、ハンズフリーで容易に記録を行える環境が強く求められている。

近年、深層学習に基づく音響モデルや大規模基盤モデルの進化により、音声認識および自然言語処理の精度が飛躍的に向上している。
これにより、現場の環境音や雑音が存在する状況であっても正確な発話テキスト化が可能となり、さらにそのテキストから作業の意図や詳細分類を抽出できる技術的基盤が整いつつある。
しかしながら、計算資源が著しく制限される小型エッジマイコン上においては、高度なノイズキャンセリングや大規模言語モデルを直接動作させることができない。そのため、突発的な環境騒音や第三者の話し声に起因する音声検知の誤作動をいかに軽量な構造で防止するかという技術的・学術的課題が存在している。

\subsection{本研究の目的と学術的な問い}
本研究の目的は、小型エッジ端末（M5CoreS3）とサーバ側での高度なAI解析（WhisperおよびGPT）を連携させ、主要な作業工程に絞った作業ログ自動化システムを構築することである。

本研究における学術的な問いは以下の通りである。
\begin{quote}
計算リソースが極めて制限されたエッジマイコン（ESP32）環境において、対象フレーズを『加工開始』『加工完了』『段取開始』『段取終了』の4種類に限定し、物理空間情報である『BLE電波強度（RSSI）』と『音声・言語特徴（Whisper/GPT-4）』をマルチモーダルに結合・制御することで、エッジ側の計算負荷を高めることなく環境騒音や第三者の発話による音声誤検知・誤記録をどこまで低減できるか？
\end{quote}

ユーザは作業中に端末のタッチ操作または発話（「加工開始」「加工完了」「段取開始」「段取終了」等）を行うのみで、音声データが即座に解析され、該当する作業状態および現場の関連情報（BLEビーコンによる位置情報等）が自動的にデータベースへ記録される。
提案手法により、エッジマイコンの軽量性を維持しつつ、現場作業者の記録負担を極力排除し、正確な作業履歴の蓄積と業務プロセスの可視化・効率化を実現する。

\subsection{本書の構成}
本書の構成は以下の通りである。
第\ref{sec2}章では、本研究の基礎となる関連技術および先行研究について述べる。
第\ref{sec3}章では、提案システムの全体構成および音声認識・データ送信・自然言語処理などの実装内容について詳述する。
第\ref{sec4}章では、構築したシステムを用いた評価実験の結果とその考察について述べる。
第\ref{sec5}章では、本研究の結論および今後の課題についてまとめる。

\newpage

\section{関連技術・関連研究}\label{sec2}
本研究に関連する技術的要素および先行研究について、「音声処理・言語モデル」と「エッジ端末・近距離通信」の観点から述べる。

\subsection{関連技術}
\subsubsection{Whisper}
Whisper \cite{radford2023whisper} は、OpenAIが開発した大規模音声認識モデルであり、多言語対応および高い雑音耐性を特徴とする。
様々な音響環境下での発話に対しても高い認識精度を誇り、本研究では入力音声データを正確なテキストへ変換する基盤処理として採用する。

\subsubsection{大規模言語モデル（GPT）}
GPT-4 \cite{openai2023gpt4} などの大規模言語モデル（LLM）は、文脈の高度な理解と構造化データ変換能力を有する。
Weiら \cite{wei2023zero} の研究が示すように、LLMに対するプロンプト制御（Zero-Shot抽出）を用いることで、音声認識テキストに含まれる語尾のゆらぎや騒音による微細な誤字を吸収し、「加工開始」「加工完了」「段取開始」「段取終了」の定義済みカテゴリへ正確にマッピングすることが可能となる。

\subsubsection{M5CoreS3およびBLE通信}
M5CoreS3はESP32-S3を搭載した小型マイクロコントローラであり、タッチスクリーン、マイク、Wi-Fi/BLE機能を標準で備えている。
Bluetooth Low Energy（BLE）通信を利用することで、周囲に設置されたビーコン機器の電波強度（RSSI）をスキャンし、作業者の滞在エリアや使用対象の機械・設備を特定する手段として活用できる。

\subsection{関連研究}
\subsubsection{音声認識の業務利用と課題}
音声入力を活用した作業記録や業務支援システムは、医療・看護記録やインフラ設備の保守点検分野を中心に研究が進められてきた。
河原 \cite{kawahara2016} は、現場作業における音声入力の実用化において、環境騒音、発話の揺らぎ、および未知語への対応が大きな障壁になると指摘している。
また、大西ら \cite{onishi2018} は保守点検業務における音声入力システムの試作を行い、作業者の手を止めない入力インターフェースの有効性を確認している。
さらに、佐藤ら \cite{sato2020} はBLEビーコンの電波強度（RSSI）を用いた位置特定手法を提案し、作業者の位置追跡の有用性を示している。

本研究はこれらの先行研究に対し、大語彙・雑音耐性を持つWhisperとプロンプト制御によるGPTを融合させることで、現場での発話を定型4フレーズ（加工開始・加工完了・段取開始・段取終了）へ正確に自動構造化する点、ならびにエッジ端末の近距離無線情報と連動させて誤入力・誤起動を抑制する点に新規性がある。

\newpage

\section{研究内容}\label{sec3}
本研究では、M5CoreS3端末を用いて音声データおよび環境情報を取得し、Wi-Fi経由でバックエンドサーバへ送信して解析・蓄積を行うシステムを構築した。

\subsection{システム全体構成}
システムの全体構成を図に定義される処理フローに基づいて構築する。本システムは大きく分けて「エッジ端末処理部」と「サーバ解析・記録部」から構成される。

\begin{itemize}
    \item \textbf{エッジ端末（M5CoreS3）}: 内蔵マイクによる音声録音（16kHz PCM）、BLEビーコン信号のスキャン、タッチUIの管理、サーバへのHTTP/TCP通信を担当。
    \item \textbf{サーバ部}: Web API（Python Flask/FastAPI等）、Whisperによる音声テキスト化処理、GPT APIによる4フレーズ判定・構造化解析、ログデータベースへの保存処理を担当。
\end{itemize}

\subsection{音声認識・データ送信処理}
エッジ端末側では、PSRAM上に録音バッファを動的に確保し、最高10秒間の音声データ（16kHz, 16bit PCM）を取得する。
画面タップにより録音が開始・停止され、取得されたバイナリデータは直ちにHTTP POSTリクエスト経由でサーバの指定エンドポイントへ転送される。

サーバ側では、受信した音声バイナリをWhisperモデルに入力し、即座にテキストに変換する。
変換結果はHTTPレスポンスとしてM5CoreS3端末へ返送され、端末の画面上にリアルタイム表示されることで、ユーザーが認識内容をその場で確認・確認保存できるUIを実装した。

\subsection{BLEによる位置・対象識別処理}
M5CoreS3はデュアルコア（Core 0 / Core 1）構成となっており、音声処理やUI描画を行うメインタスク（Core 1）とは独立して、Core 0 上でバックグラウンドでのBLEスキャンタスクを実行する。
周囲に存在するターゲットビーコンのRSSI値を継続的に取得・平均化することで、作業者が現在どの機械・作業エリアに近接しているかを自動判定し、音声ログ送信時にその位置・対象メタデータを自動で付与する仕組みを構築した。

\subsection{自然言語処理（GPT）による意図解析とログ構造化}
Whisperによって得られた発話テキストは、サーバ上の言語処理モジュールに渡される。
プロンプトエンジニアリングを活用し、認識テキストを以下の4つの定義済みフレーズ（および例外）に限定して分類・変換を行う。

\begin{itemize}
    \item \textbf{加工開始}: 加工工程の開始を識別
    \item \textbf{加工完了}: 加工工程の終了を識別
    \item \textbf{段取開始}: 段取り・準備工程の開始を識別
    \item \textbf{段取終了}: 段取り・準備工程の終了を識別
\end{itemize}

解析結果は以下の要素を含むJSON構造データへと自動変換され、データベースへ記録される。
\begin{itemize}
    \item \textbf{Action}: 定義された作業種別（加工開始 / 加工完了 / 段取開始 / 段取終了）
    \item \textbf{Target\_Area}: BLEから判定された作業エリア・対象機械ID
    \item \textbf{Timestamp}: 記録日時（NTPにより同期された正確な時刻）
\end{itemize}

\subsection{利用ツール・開発環境}
本システムの開発および検証には以下のツール・環境を利用した。
\begin{itemize}
    \item \textbf{ハードウェア}: M5CoreS3 (ESP32-S3)
    \item \textbf{開発言語・環境}: C++ (Arduino IDE / PlatformIO), Python 3.10
    \item \textbf{使用AIモデル}: OpenAI Whisper, OpenAI GPT-4 API
    \item \textbf{通信プロトコル}: HTTP/1.1, TCP/IP, BLE (Bluetooth Low Energy)
\end{itemize}

\newpage

\section{結果・考察}\label{sec4}

\subsection{実験および検証}
構築したシステムを用いて、静寂な室内環境および動作音が存在する作業現場環境（近畿工業株式会社の現場環境等）を模した条件化において、発話・録音・解析・保存の一連の動作検証を実施した。

実証実験では、作業者が「加工開始」「加工完了」「段取開始」「段取終了」の4フレーズ、およびそれに類する口語表現（例: 「加工を始めます」「段取り終わりました」）を M5CoreS3 に向けて発声し、ログ化の成功率および処理にかかるレスポンス時間を測定した。

\subsection{評価結果}
実験の結果、以下の知見が得られた。

\begin{enumerate}
    \item \textbf{認識精度と定型分類の正確性}: 対象を4フレーズに限定しWhisperとGPTを組み合わせることで、「加工開始」「段取終了」等の定型フレーズだけでなく、「加工を始める」「段取りが終わった」といった同義の口語発話からも、高精度に目的の4カテゴリへ正確に分類できることが確認された。
    \item \textbf{マルチコア並行処理の安定性}: Core 0 でBLEスキャンを行い、Core 1 でマイク録音・表示処理を行うマルチタスク設計により、音声処理中にBLEスキャン処理が干渉することなく、円滑な通信と画面応答が維持された。
    \item \textbf{応答時間}: 音声録音終了からサーバ側でのWhisper/GPT解析を経て端末画面に結果がフィードバックされるまでの平均レスポンス時間は約2.5秒～4.0秒であり、実用上許容できる範囲であることが確認された。
\end{enumerate}

\subsection{考察}
現場環境における実運用においては、周囲の騒音レベルが高い場合にWhisperの誤認識率が僅かに上昇する傾向が見られた。
しかし、GPT側で対象を「加工開始」「加工完了」「段取開始」「段取終了」の4候補に絞り込むプロンプト条件を付与したことで、音声認識結果に一部誤字が含まれる場合であっても、正しい作業カテゴリへ正しく補正・分類され、高い堅牢性が示された。

また、M5CoreS3のタッチスクリーン操作における誤検知（静電容量タッチのノイズ等）に対しては、ソフトウェア側でガードタイマーおよび確実なタップ（`wasClicked`）判定条件を導入したことで、意図しない自動録音のループ挙動を完全に抑制できた。

\newpage

\section{結論・今後の課題}\label{sec5}

\subsection{結論}
本研究では、M5CoreS3デバイスとWhisper・GPT等の先進的AIモデルを融合させたハンズフリー作業ログ自動記録システムを構築・評価した。
本研究で設定した学術的問いに対し、認識フレーズを「加工開始」「加工完了」「段取開始」「段取終了」の4つに集約し、BLE電波強度による物理空間制約とWhisper/GPTによる意味解析を統合することで、リソースの限られたエッジデバイス上であっても環境騒音の影響を受けにくい高精度な作業ログ自動化が実現できることを示した。

本システムにより、以下の成果が得られた。
\begin{itemize}
    \item キーボード等の入力操作を必要としない、音声を中心とした作業記録環境の実現。
    \item 「加工開始」「加工完了」「段取開始」「段取終了」の4フレーズに基づく作業状態およびBLEによる位置情報の自動構造化。
    \item アナログな記録作業に伴う作業中断や入力コストの劇的な削減。
\end{itemize}

\subsection{今後の課題}
今後の展望および改良課題として以下が挙げられる。
\begin{itemize}
    \item \textbf{超極小遅延化}: エッジ側での局所的な軽量音声認識（Local Whisper等）の導入によるレスポンス時間のさらなる短縮。
    \item \textbf{騒音適応性の向上}: ノイズキャンセリングマイクの採用や、現場特有の機械音に対する音響フィルタ処理の追加。
    \item \textbf{大規模現場での実証}: 近畿工業株式会社などの実際の工場ラインにおける長期間の運用試験と、作業効率化に関する定量的データの蓄積。
\end{itemize}

\newpage

\section*{謝辞}
\addcontentsline{toc}{section}{謝辞}
本研究に取り組むにあたり、親身なご指導と温かいご鞭撻を賜りました指導教員の森山真光准教授に心より感謝申し上げます。
また、実証実験の機会や貴重な現場データをご提供いただきました近畿工業株式会社の皆様に深く御礼申し上げます。
最後に、日々活発な議論を行い、実験環境の構築にご協力いただきました研究室の皆様に感謝いたします。

\newpage

\begin{thebibliography}{99}
\addcontentsline{toc}{section}{参考文献}
\bibitem{radford2023whisper}
Alec Radford, Jong Wook Kim, Tao Xu, Greg Brockman, Christine McLeavey, and Ilya Sutskever.
Robust Speech Recognition via Large-Scale Weak Supervision.
In {\it Proceedings of the 40th International Conference on Machine Learning (ICML)}, pp. 28492--28518, 2023.

\bibitem{openai2023gpt4}
OpenAI.
GPT-4 Technical Report.
arXiv preprint arXiv:2303.08774, 2023.

\bibitem{wei2023zero}
Jason Wei, Xuezhi Wang, Dale Schuurmans, Maarten Bosma, Brian Ichter, Fei Xia, Ed Chi, Quoc Le, and Denny Zhou.
Chain-of-Thought Prompting Elicits Reasoning in Large Language Models.
In {\it Advances in Neural Information Processing Systems (NeurIPS)}, Vol. 35, pp. 24824--24837, 2022.

\bibitem{kawahara2016}
河原 英紀.
音声認識技術の現状と現場業務支援への応用課題.
日本音響学会誌, Vol. 72, No. 8, pp. 485--492, 2016.

\bibitem{onishi2018}
大西 智和, 他.
音声認識技術を活用した現場作業支援システムの試作と評価.
情報処理学会研究報告 (MPS), Vol. 2018-MPS-117, No. 10, pp. 1--6, 2018.

\bibitem{sato2020}
佐藤 裕樹, 他.
BLEビーコンを用いた位置推定技術の現場作業管理への適用.
電気学会論文誌Ｃ（電子・情報・システム部門誌）, Vol. 140, No. 5, pp. 512--519, 2020.
\end{thebibliography}

\newpage
\appendix

\end{document}